
\documentclass{article}
\usepackage{colortbl}
\usepackage{makecell}
\usepackage{multirow}
\usepackage{supertabular}

\begin{document}

\newcounter{utterance}

\centering \large Interaction Transcript for game `dond', experiment `coop', episode 1 with gpt{-}5.4{-}azure.
\vspace{24pt}

{ \footnotesize  \setcounter{utterance}{1}
\setlength{\tabcolsep}{0pt}
\begin{supertabular}{c@{$\;$}|p{.15\linewidth}@{}p{.15\linewidth}p{.15\linewidth}p{.15\linewidth}p{.15\linewidth}p{.15\linewidth}}
    \# & $\;$A & \multicolumn{4}{c}{Game Master} & $\;\:$B\\
    \hline

    \theutterance \stepcounter{utterance}  
    & & \multicolumn{4}{p{0.6\linewidth}}{
        \cellcolor[rgb]{0.9,0.9,0.9}{
            \makecell[{{p{\linewidth}}}]{
                \texttt{\tiny{[P1$\langle$GM]}}
                \texttt{Du spiller et kooperativt forhandlingsspil, hvor du skal enes med en anden spiller om, hvordan du deler et sæt genstande mellem jer.} \\
\\ 
\texttt{Regler:} \\
\texttt{(a) Du og den anden spiller får et sæt genstande. Du får også en hemlig værdifunktion, der viser, hvor meget du værdisætter hver type objekt.} \\
\texttt{(b) Du bytter beskeder med den anden spiller for at blive enig om, hvem der får hvad. Du kan sende maksimalt 5 beskeder hver eller afslutte spillet tidligt ved at fremsætte et hemmeligt forslag når som helst.} \\
\texttt{(c) Du skal indsende et hemmeligt forslag, der angiver de genstande, du vil have, formateret i firkantede parenteser som følger: "{[}Forslag: <antal> <objektnavn>, <antal> <objektnavn>, <...>{]}"} \\
\texttt{(d) Hvis dine forslag er komplementære, dvs. der er nok genstande til at opfylde begge forslag, tildeles hver spiller en score baseret på summen af værdierne for de genstande, hun modtog. Ellers får du begge nul point.} \\
\texttt{(e) Dit mål er at maksimere summen af din score og den anden spillers score.} \\
\\ 
\texttt{Lad os starte.} \\
\\ 
\texttt{Sættet af tilgængelige genstande er:} \\
\\ 
\texttt{2 penne, 2 citroner, 2 potter} \\
\\ 
\texttt{Din hemmelige værdifunktion er:} \\
\\ 
\texttt{pen: 2, citron: 0, potte: 3} \\
\\ 
\texttt{VIGTIG: Dine beskeder, medmindre det er et hemmeligt forslag, bliver sendt direkte til den anden spiller, så inkluder ikke noget svar på reglerne eller tekst, der annoncerer din besked. Brug det angivne format for at afgive et hemmeligt forslag. Brug ikke firkantede parenteser, når du kommunikerer med den anden spiller, eller det vil blive fortolket som dit hemmelige forslag.} \\
            }
        }
    }
    & \\ \\

    \theutterance \stepcounter{utterance}  
    & \multicolumn{4}{p{0.6\linewidth}}{
        \cellcolor[rgb]{0.9,0.9,0.9}{
            \makecell[{{p{\linewidth}}}]{
                \texttt{\tiny{[P1$\rangle$GM]}}
                \texttt{Hvilke genstande er vigtigst for dig? Citroner er lav prioritet for mig, mens penne og især potter betyder mere. Hvis du siger, hvad du helst vil have, kan vi lave en effektiv deling.} \\
            }
        }
    }
    & & \\ \\

    \theutterance \stepcounter{utterance}  
    & & & \multicolumn{2}{p{0.3\linewidth}}{
        \cellcolor[rgb]{0.9,0.9,0.9}{
            \makecell[{{p{\linewidth}}}]{
                \texttt{\tiny{[GM$|$GM]}}
                \texttt{continue} \\
            }
        }
    }
    & & \\ \\

    \theutterance \stepcounter{utterance}  
    & & & \multicolumn{2}{p{0.3\linewidth}}{
        \cellcolor[rgb]{0.9,0.9,0.9}{
            \makecell[{{p{\linewidth}}}]{
                \texttt{\tiny{[GM$|$GM]}}
                \texttt{Hvilke genstande er vigtigst for dig? Citroner er lav prioritet for mig, mens penne og især potter betyder mere. Hvis du siger, hvad du helst vil have, kan vi lave en effektiv deling.} \\
            }
        }
    }
    & & \\ \\

    \theutterance \stepcounter{utterance}  
    & & \multicolumn{4}{p{0.6\linewidth}}{
        \cellcolor[rgb]{0.9,0.9,0.9}{
            \makecell[{{p{\linewidth}}}]{
                \texttt{\tiny{[GM$\rangle$P2]}}
                \texttt{Du spiller et kooperativt forhandlingsspil, hvor du skal enes med en anden spiller om, hvordan du deler et sæt genstande mellem jer.} \\
\\ 
\texttt{Regler:} \\
\texttt{(a) Du og den anden spiller får et sæt genstande. Du får også en hemlig værdifunktion, der viser, hvor meget du værdisætter hver type objekt.} \\
\texttt{(b) Du bytter beskeder med den anden spiller for at blive enig om, hvem der får hvad. Du kan sende maksimalt 5 beskeder hver eller afslutte spillet tidligt ved at fremsætte et hemmeligt forslag når som helst.} \\
\texttt{(c) Du skal indsende et hemmeligt forslag, der angiver de genstande, du vil have, formateret i firkantede parenteser som følger: "{[}Forslag: <antal> <objektnavn>, <antal> <objektnavn>, <...>{]}"} \\
\texttt{(d) Hvis dine forslag er komplementære, dvs. der er nok genstande til at opfylde begge forslag, tildeles hver spiller en score baseret på summen af værdierne for de genstande, hun modtog. Ellers får du begge nul point.} \\
\texttt{(e) Dit mål er at maksimere summen af din score og den anden spillers score.} \\
\\ 
\texttt{Lad os starte.} \\
\\ 
\texttt{Sættet af tilgængelige genstande er:} \\
\\ 
\texttt{2 penne, 2 citroner, 2 potter} \\
\\ 
\texttt{Din hemmelige værdifunktion er:} \\
\\ 
\texttt{pen: 3, citron: 1, potte: 1} \\
\\ 
\texttt{VIGTIG: Dine beskeder, medmindre det er et hemmeligt forslag, bliver sendt direkte til den anden spiller, så inkluder ikke noget svar på reglerne eller tekst, der annoncerer din besked. Brug det angivne format for at afgive et hemmeligt forslag. Brug ikke firkantede parenteser, når du kommunikerer med den anden spiller, eller det vil blive fortolket som dit hemmelige forslag.} \\
\\ 
\texttt{Hvilke genstande er vigtigst for dig? Citroner er lav prioritet for mig, mens penne og især potter betyder mere. Hvis du siger, hvad du helst vil have, kan vi lave en effektiv deling.} \\
            }
        }
    }
    & \\ \\

    \theutterance \stepcounter{utterance}  
    & & & \multicolumn{4}{p{0.6\linewidth}}{
        \cellcolor[rgb]{0.9,0.9,0.9}{
            \makecell[{{p{\linewidth}}}]{
                \texttt{\tiny{[GM$\langle$P2]}}
                \texttt{Penne er klart vigtigst for mig; citroner og potter er ret lave for mig. Hvis potter er vigtigst for dig, virker en effektiv deling som: jeg tager 2 penne, og du tager 2 potter. Du kan også få citronerne, hvis de ikke betyder meget for mig.} \\
            }
        }
    }
     \\ \\

    \theutterance \stepcounter{utterance}  
    & & & \multicolumn{2}{p{0.3\linewidth}}{
        \cellcolor[rgb]{0.9,0.9,0.9}{
            \makecell[{{p{\linewidth}}}]{
                \texttt{\tiny{[GM$|$GM]}}
                \texttt{continue} \\
            }
        }
    }
    & & \\ \\

    \theutterance \stepcounter{utterance}  
    & & & \multicolumn{2}{p{0.3\linewidth}}{
        \cellcolor[rgb]{0.9,0.9,0.9}{
            \makecell[{{p{\linewidth}}}]{
                \texttt{\tiny{[GM$|$GM]}}
                \texttt{Penne er klart vigtigst for mig; citroner og potter er ret lave for mig. Hvis potter er vigtigst for dig, virker en effektiv deling som: jeg tager 2 penne, og du tager 2 potter. Du kan også få citronerne, hvis de ikke betyder meget for mig.} \\
            }
        }
    }
    & & \\ \\

    \theutterance \stepcounter{utterance}  
    & & \multicolumn{4}{p{0.6\linewidth}}{
        \cellcolor[rgb]{0.9,0.9,0.9}{
            \makecell[{{p{\linewidth}}}]{
                \texttt{\tiny{[P1$\langle$GM]}}
                \texttt{Penne er klart vigtigst for mig; citroner og potter er ret lave for mig. Hvis potter er vigtigst for dig, virker en effektiv deling som: jeg tager 2 penne, og du tager 2 potter. Du kan også få citronerne, hvis de ikke betyder meget for mig.} \\
            }
        }
    }
    & \\ \\

    \theutterance \stepcounter{utterance}  
    & \multicolumn{4}{p{0.6\linewidth}}{
        \cellcolor[rgb]{0.9,0.9,0.9}{
            \makecell[{{p{\linewidth}}}]{
                \texttt{\tiny{[P1$\rangle$GM]}}
                \texttt{Det lyder som en god og effektiv deling: du tager 2 penne, og jeg tager 2 potter. Du må også gerne tage citronerne, for de er ikke vigtige for mig.} \\
            }
        }
    }
    & & \\ \\

    \theutterance \stepcounter{utterance}  
    & & & \multicolumn{2}{p{0.3\linewidth}}{
        \cellcolor[rgb]{0.9,0.9,0.9}{
            \makecell[{{p{\linewidth}}}]{
                \texttt{\tiny{[GM$|$GM]}}
                \texttt{continue} \\
            }
        }
    }
    & & \\ \\

    \theutterance \stepcounter{utterance}  
    & & & \multicolumn{2}{p{0.3\linewidth}}{
        \cellcolor[rgb]{0.9,0.9,0.9}{
            \makecell[{{p{\linewidth}}}]{
                \texttt{\tiny{[GM$|$GM]}}
                \texttt{Det lyder som en god og effektiv deling: du tager 2 penne, og jeg tager 2 potter. Du må også gerne tage citronerne, for de er ikke vigtige for mig.} \\
            }
        }
    }
    & & \\ \\

    \theutterance \stepcounter{utterance}  
    & & \multicolumn{4}{p{0.6\linewidth}}{
        \cellcolor[rgb]{0.9,0.9,0.9}{
            \makecell[{{p{\linewidth}}}]{
                \texttt{\tiny{[GM$\rangle$P2]}}
                \texttt{Det lyder som en god og effektiv deling: du tager 2 penne, og jeg tager 2 potter. Du må også gerne tage citronerne, for de er ikke vigtige for mig.} \\
            }
        }
    }
    & \\ \\

    \theutterance \stepcounter{utterance}  
    & & & \multicolumn{4}{p{0.6\linewidth}}{
        \cellcolor[rgb]{0.9,0.9,0.9}{
            \makecell[{{p{\linewidth}}}]{
                \texttt{\tiny{[GM$\langle$P2]}}
                \texttt{{[}Forslag: 2 penne, 2 citroner{]}} \\
            }
        }
    }
     \\ \\

    \theutterance \stepcounter{utterance}  
    & & & \multicolumn{2}{p{0.3\linewidth}}{
        \cellcolor[rgb]{0.9,0.9,0.9}{
            \makecell[{{p{\linewidth}}}]{
                \texttt{\tiny{[GM$|$GM]}}
                \texttt{proposal} \\
            }
        }
    }
    & & \\ \\

    \theutterance \stepcounter{utterance}  
    & & & \multicolumn{2}{p{0.3\linewidth}}{
        \cellcolor[rgb]{0.9,0.9,0.9}{
            \makecell[{{p{\linewidth}}}]{
                \texttt{\tiny{[GM$|$GM]}}
                [2, 2, 0]
            }
        }
    }
    & & \\ \\

    \theutterance \stepcounter{utterance}  
    & & \multicolumn{4}{p{0.6\linewidth}}{
        \cellcolor[rgb]{0.9,0.9,0.9}{
            \makecell[{{p{\linewidth}}}]{
                \texttt{\tiny{[P1$\langle$GM]}}
                \texttt{Den anden spiller har indsendt et hemmeligt forslag. Nu er det din tur til at indsende dit eget forslag.} \\
            }
        }
    }
    & \\ \\

    \theutterance \stepcounter{utterance}  
    & \multicolumn{4}{p{0.6\linewidth}}{
        \cellcolor[rgb]{0.9,0.9,0.9}{
            \makecell[{{p{\linewidth}}}]{
                \texttt{\tiny{[P1$\rangle$GM]}}
                \texttt{{[}Forslag: 2 potter{]}} \\
            }
        }
    }
    & & \\ \\

    \theutterance \stepcounter{utterance}  
    & & & \multicolumn{2}{p{0.3\linewidth}}{
        \cellcolor[rgb]{0.9,0.9,0.9}{
            \makecell[{{p{\linewidth}}}]{
                \texttt{\tiny{[GM$|$GM]}}
                \texttt{proposal} \\
            }
        }
    }
    & & \\ \\

    \theutterance \stepcounter{utterance}  
    & & & \multicolumn{2}{p{0.3\linewidth}}{
        \cellcolor[rgb]{0.9,0.9,0.9}{
            \makecell[{{p{\linewidth}}}]{
                \texttt{\tiny{[GM$|$GM]}}
                [0, 0, 2]
            }
        }
    }
    & & \\ \\

    \theutterance \stepcounter{utterance}  
    & & & \multicolumn{2}{p{0.3\linewidth}}{
        \cellcolor[rgb]{0.9,0.9,0.9}{
            \makecell[{{p{\linewidth}}}]{
                \texttt{\tiny{[GM$|$GM]}}
                [[0, 0, 2], [2, 2, 0]]
            }
        }
    }
    & & \\ \\

\end{supertabular}
}

\end{document}
